%% toc: P2 14 14.7
\section{Proofs for interpolation at the boundary}
\label{sec:annex-p02-w02-p02-015-interpolation-boundary}

Throughout, \(\varphi\entails\psi\),
\(X=\operatorname{Var}(\varphi)\cap\operatorname{Var}(\psi)\) is listed as
\(x_1,\ldots,x_k\), and \(\kappa_u\), \(U_\varphi=\{u_1,\ldots,u_m\}\) and
the output \(\theta\) are as in
Definition~\ref{def:p02-w02-p02-015-output}.  The constant convention is
Convention~\ref{def:p02-w02-p02-015-constants}, and ``coincidence'' refers to
Lemma~\ref{lem:prop-coincidence}.

\begin{annexitem}{annex:p02-w02-p02-015-content}{thm:p02-w02-p02-015-output-content}{Semantic content of the output}
\begin{proof}
\emph{Clause 1.}  Suppose first \(k=0\).  Then \(\kappa_u\) is the empty
conjunction \(\topv\), that is \(\neg\bottom\).  Since
\(\widehat v(\bottom)=\mathbf{F}\), the negation clause gives
\(\widehat v(\neg\bottom)=\mathbf{T}\), so \(v\entails\kappa_u\).  On the other
side, \(u\) and \(v\mathbin{\upharpoonright}\varnothing\) are both the empty
function \(u_\varnothing\), so \(v\mathbin{\upharpoonright}X=u\).  Both sides
of the equivalence hold.

Suppose \(k\geq 1\).  Let \(\ell_i\) be the literal for \(u\) at \(x_i\), put
\(C_k=\ell_k\) and \(C_i=(\ell_i\land C_{i+1})\) for \(1\leq i<k\), so that
\(\kappa_u=C_1\).  First, \(v\entails\ell_i\) if and only if
\(v(x_i)=u(x_i)\): if \(u(x_i)=1\), then \(\ell_i=x_i\), and
\(v\entails x_i\) iff \(v(x_i)=\mathbf{T}\); if \(u(x_i)=0\), then
\(\ell_i=\neg x_i\), and by the negation clause \(v\entails\neg x_i\) iff
\(v(x_i)=\mathbf{F}\).  Second, by downward induction on \(i\),
\(v\entails C_i\) if and only if \(v\entails\ell_j\) for every \(j\) with
\(i\leq j\leq k\).  For \(i=k\) this is the equation \(C_k=\ell_k\).  For
\(i<k\) the conjunction clause gives \(v\entails(\ell_i\land C_{i+1})\) iff
\(v\entails\ell_i\) and \(v\entails C_{i+1}\), and the induction hypothesis
rewrites the second conjunct.  Taking \(i=1\) and combining the two facts,
\(v\entails\kappa_u\) iff \(v(x_j)=u(x_j)\) for \(j=1,\ldots,k\), that is iff
\(v\mathbin{\upharpoonright}X=u\).

\emph{Clause 2.}  If \(m=0\), then \(\theta=\bottom\); no valuation satisfies
\(\bottom\), and \(U_\varphi=\varnothing\) contains no
\(v\mathbin{\upharpoonright}X\); both sides fail for every \(v\).  If
\(m\geq 1\), put \(D_m=\kappa_{u_m}\) and
\(D_j=(\kappa_{u_j}\lor D_{j+1})\) for \(1\leq j<m\), so that
\(\theta=D_1\).  By downward induction on \(j\), using the disjunction clause
at each step, \(v\entails D_j\) iff \(v\entails\kappa_{u_l}\) for some \(l\)
with \(j\leq l\leq m\).  With \(j=1\) and Clause~1, \(v\entails\theta\) iff
\(v\mathbin{\upharpoonright}X=u_l\) for some \(l\), that is iff
\(v\mathbin{\upharpoonright}X\in U_\varphi\).

\emph{Clause 3.}  If \(k=0\), \(\operatorname{Var}(\kappa_u)=
\operatorname{Var}(\topv)=\varnothing\).  If \(k\geq 1\),
\(\operatorname{Var}(\ell_i)=\{x_i\}\) by the variable and negation clauses
of the occurrence set, and the conjunction clause gives
\(\operatorname{Var}(\kappa_u)=\{x_1,\ldots,x_k\}=X\).  If \(m=0\),
\(\operatorname{Var}(\theta)=\operatorname{Var}(\bottom)=\varnothing\).  If
\(m\geq 1\), the disjunction clause gives
\(\operatorname{Var}(\theta)=\bigcup_{j=1}^{m}\operatorname{Var}(\kappa_{u_j})\),
which is \(\varnothing\) when \(k=0\) and \(X\) when \(k\geq 1\).  In all cases
\(\operatorname{Var}(\theta)\subseteq X\), with equality to \(\varnothing\)
when \(X=\varnothing\) or \(U_\varphi=\varnothing\).
\end{proof}
\end{annexitem}

\begin{annexitem}{annex:p02-w02-p02-015-varfree}{thm:p02-w02-p02-015-variable-free}{Variable-free formulas are constant}
\begin{proof}
Let \(\operatorname{Var}(\chi)=\varnothing\) and let \(v,w\) be valuations.
No variable occurs in \(\chi\), so \(v\) and \(w\) agree on every variable
occurring in \(\chi\), the condition having no instance.  Coincidence gives
\(\widehat v(\chi)=\widehat w(\chi)\).  Hence all valuations give \(\chi\) the
same truth value: if that value is \(\mathbf{T}\), every valuation satisfies
\(\chi\); if it is \(\mathbf{F}\), none does.  For \(\bottom\) the valuation clause
gives \(\widehat v(\bottom)=\mathbf{F}\), and for \(\topv=\neg\bottom\) the
negation clause then gives \(\widehat v(\topv)=\mathbf{T}\), for every \(v\).
\end{proof}
\end{annexitem}

\begin{annexitem}{annex:p02-w02-p02-015-empty}{thm:p02-w02-p02-015-empty-shared}{Empty shared vocabulary}
\begin{proof}
\emph{Step 1: the selected family.}  Since \(X=\varnothing\), the only
assignment \(X\to\{0,1\}\) is \(u_\varnothing\), and
\(v\mathbin{\upharpoonright}\varnothing=u_\varnothing\) for every valuation
\(v\).  Hence \(u_\varnothing\in U_\varphi\) iff some valuation satisfies
\(\varphi\).  So \(U_\varphi=\{u_\varnothing\}\) when \(\varphi\) is
satisfiable and \(U_\varphi=\varnothing\) when it is not.

\emph{Step 2: the complete conjunction.}  With \(k=0\) the complete
conjunction \(\kappa_{u_\varnothing}\) is the empty conjunction, which by the
convention is \(\topv\).

\emph{Clause 1.}  Let \(\varphi\) be satisfiable.  Then \(m=1\) and \(\theta\)
is the disjunction of the one-term sequence \(\kappa_{u_\varnothing}\), which
by the convention for \(n=1\) is \(\kappa_{u_\varnothing}=\topv\).

\(\varphi\entails\topv\): every valuation satisfies \(\topv\)
(Lemma~\ref{thm:p02-w02-p02-015-variable-free}), in particular every valuation
satisfying \(\varphi\).

\(\psi\) is valid: let \(v\) be any valuation, and fix a valuation \(w\)
satisfying \(\varphi\).  Define \(v'\) by \(v'(p)=w(p)\) when
\(p\in\operatorname{Var}(\varphi)\) and \(v'(p)=v(p)\) otherwise.  Then \(v'\)
agrees with \(w\) on \(\operatorname{Var}(\varphi)\), so by coincidence
\(v'\entails\varphi\), and the hypothesis \(\varphi\entails\psi\) gives
\(v'\entails\psi\).  If \(p\in\operatorname{Var}(\psi)\), then
\(p\notin\operatorname{Var}(\varphi)\) because
\(\operatorname{Var}(\varphi)\cap\operatorname{Var}(\psi)=\varnothing\), so
\(v'(p)=v(p)\).  Thus \(v'\) and \(v\) agree on
\(\operatorname{Var}(\psi)\), and coincidence gives \(v\entails\psi\).  As
\(v\) was arbitrary, \(\psi\) is valid.

\(\topv\entails\psi\): every valuation satisfies \(\psi\), in particular every
valuation satisfying \(\topv\).

\emph{Clause 2.}  Let \(\varphi\) be unsatisfiable.  Then \(U_\varphi=
\varnothing\), \(m=0\), and \(\theta\) is the empty disjunction, which by the
convention is \(\bottom\).

\(\varphi\entails\bottom\): the requirement that every valuation satisfying
\(\varphi\) satisfy \(\bottom\) has no instance, because no valuation
satisfies \(\varphi\).

\(\bottom\entails\psi\): no valuation satisfies \(\bottom\)
(Lemma~\ref{thm:p02-w02-p02-015-variable-free}), so the requirement that every
valuation satisfying \(\bottom\) satisfy \(\psi\) has no instance.

\emph{Clause 3.}  By the occurrence-set clauses,
\(\operatorname{Var}(\bottom)=\varnothing\) and
\(\operatorname{Var}(\topv)=\operatorname{Var}(\neg\bottom)=
\operatorname{Var}(\bottom)=\varnothing\).  In either case
\(\operatorname{Var}(\theta)=\varnothing\subseteq\operatorname{Var}(\varphi)
\cap\operatorname{Var}(\psi)\), so the statement that every variable of
\(\theta\) occurs in both \(\varphi\) and \(\psi\) holds, having no instance.
With Clauses~1 and~2, \(\theta\) meets all three conditions of
Definition~\ref{def:p02-interpolant}.

The argument used the variable-free formulas \(\topv\) and \(\bottom\) at
Step~2 and Clause~2; that these exist in \(\mathsf{Fm}_{\mathsf{CPL}}\) is the formation clause
for \(\bottom\) together with the abbreviation \(\topv=\neg\bottom\).
\end{proof}
\end{annexitem}

\begin{annexitem}{annex:p02-w02-p02-015-dichot}{thm:p02-w02-p02-015-dichotomy}{Consequence without shared variables}
\begin{proof}
If \(\varphi\) is unsatisfiable there is nothing to prove.  If \(\varphi\) is
satisfiable, Clause~1 of Theorem~\ref{thm:p02-w02-p02-015-empty-shared},
whose hypotheses are \(\varphi\entails\psi\) and
\(\operatorname{Var}(\varphi)\cap\operatorname{Var}(\psi)=\varnothing\), gives
that \(\psi\) is valid.
\end{proof}
\end{annexitem}

\begin{annexitem}{annex:p02-w02-p02-015-bottom}{thm:p02-w02-p02-015-constant-needed}{Need for the constant \(\bottom\)}
\begin{proof}
\emph{The description in Definition~\ref{def:p02-w02-p02-015-bottom-free}.}
\(\mathsf{Fm}_{\mathsf{CPL}}\) contains every \(p_n\) and is closed under \(\neg\) and the
three binary constructors, so leastness of \(\mathsf{Fm}_{\mathsf{CPL}}^{-}\) gives
\(\mathsf{Fm}_{\mathsf{CPL}}^{-}\subseteq\mathsf{Fm}_{\mathsf{CPL}}\).  The strings of \(\mathsf{Fm}_{\mathsf{CPL}}^{-}\) in
which \(\bottom\) does not occur contain every \(p_n\) and are closed under the
same constructors, because none of them introduces \(\bottom\); leastness
again shows that \(\bottom\) occurs in no member of \(\mathsf{Fm}_{\mathsf{CPL}}^{-}\).
Conversely, by structural induction on \(\mathsf{Fm}_{\mathsf{CPL}}\), every formula either
contains \(\bottom\) or belongs to \(\mathsf{Fm}_{\mathsf{CPL}}^{-}\): \(p_n\in
\mathsf{Fm}_{\mathsf{CPL}}^{-}\); \(\bottom\) contains \(\bottom\); if \(\chi\) contains
\(\bottom\) so does \(\neg\chi\), and otherwise \(\chi\in\mathsf{Fm}_{\mathsf{CPL}}^{-}\) and
so \(\neg\chi\in\mathsf{Fm}_{\mathsf{CPL}}^{-}\); and \((\chi\circ\chi')\) contains \(\bottom\)
if either component does, and otherwise both components lie in
\(\mathsf{Fm}_{\mathsf{CPL}}^{-}\) and so does \((\chi\circ\chi')\).  Since \(\topv\) is the
string \(\neg\bottom\), it does not occur in any member of
\(\mathsf{Fm}_{\mathsf{CPL}}^{-}\) either.

\emph{Clause 1.}  Let \(S\) be the set of \(\chi\in\mathsf{Fm}_{\mathsf{CPL}}^{-}\) with
\(\operatorname{Var}(\chi)\neq\varnothing\).  It contains every \(p_n\),
since \(\operatorname{Var}(p_n)=\{p_n\}\).  If \(\chi\in S\), then
\(\operatorname{Var}(\neg\chi)=\operatorname{Var}(\chi)\neq\varnothing\), so
\(\neg\chi\in S\).  If \(\chi,\chi'\in S\), then
\(\operatorname{Var}((\chi\circ\chi'))=\operatorname{Var}(\chi)\cup
\operatorname{Var}(\chi')\neq\varnothing\), so \((\chi\circ\chi')\in S\).
By leastness \(\mathsf{Fm}_{\mathsf{CPL}}^{-}\subseteq S\).

\emph{Clause 2.}  Let \(\chi\in\mathsf{Fm}_{\mathsf{CPL}}\) with
\(\operatorname{Var}(\chi)=\varnothing\).  By Clause~1, \(\chi\notin
\mathsf{Fm}_{\mathsf{CPL}}^{-}\), so by the converse inclusion above \(\chi\) contains
\(\bottom\).

\emph{Clause 3.}  For every valuation \(v\),
\(\widehat v(p_1\to p_1)=\mathbf{F}\) would require \(\widehat v(p_1)=\mathbf{T}\) and
\(\widehat v(p_1)=\mathbf{F}\), which is impossible; so \((p_1\to p_1)\) is valid
and in particular \(p_0\entails(p_1\to p_1)\).  Both formulas lie in
\(\mathsf{Fm}_{\mathsf{CPL}}^{-}\), the second by the binary clause applied to \(p_1\) twice,
and \(\operatorname{Var}(p_0)\cap\operatorname{Var}(p_1\to p_1)=
\{p_0\}\cap\{p_1\}=\varnothing\).  A Craig interpolant \(\chi\) for this
consequence must have every variable in this empty intersection, so
\(\operatorname{Var}(\chi)=\varnothing\), and by Clause~1 no
\(\chi\in\mathsf{Fm}_{\mathsf{CPL}}^{-}\) qualifies.

For the second consequence, \(v\entails(p_0\land\neg p_0)\) would require
\(v(p_0)=\mathbf{T}\) and, by the negation clause, \(v(p_0)=\mathbf{F}\); no valuation
satisfies the antecedent, so \((p_0\land\neg p_0)\entails p_1\) holds with no
instance to check.  Both formulas lie in \(\mathsf{Fm}_{\mathsf{CPL}}^{-}\), their occurrence
sets \(\{p_0\}\) and \(\{p_1\}\) are disjoint, and the argument of the
previous paragraph excludes every candidate in \(\mathsf{Fm}_{\mathsf{CPL}}^{-}\).

In \(\mathsf{Fm}_{\mathsf{CPL}}\), the antecedent \(p_0\) is satisfied by the valuation
sending every variable to \(\mathbf{T}\), so Clause~1 of
Theorem~\ref{thm:p02-w02-p02-015-empty-shared} gives the interpolant
\(\topv\); the antecedent \((p_0\land\neg p_0)\) is unsatisfiable, so Clause~2
of that theorem gives \(\bottom\).
\end{proof}
\end{annexitem}

\begin{annexitem}{annex:p02-w02-p02-015-incons}{thm:p02-w02-p02-015-inconsistent}{The inconsistent endpoint}
\begin{proof}
\emph{The consequence.}  No valuation satisfies \(\varphi\), so the
requirement that every valuation satisfying \(\varphi\) satisfy \(\psi\) has
no instance, and \(\varphi\entails\psi\).

\emph{The output.}  An assignment \(u:X\to\{0,1\}\) belongs to
\(U_\varphi\) only if some valuation satisfies \(\varphi\).  None does, so
\(U_\varphi=\varnothing\) for every \(X\), whether empty or not.  Hence
\(m=0\), and the output is the empty disjunction, which by the convention is
\(\bottom\).

\emph{The two entailments.}  \(\varphi\entails\bottom\) holds because no
valuation satisfies \(\varphi\).  \(\bottom\entails\psi\) holds because no
valuation satisfies \(\bottom\)
(Lemma~\ref{thm:p02-w02-p02-015-variable-free}).

\emph{The vocabulary clause.}  \(\operatorname{Var}(\bottom)=\varnothing
\subseteq X\), so every variable of \(\bottom\) occurs in both \(\varphi\) and
\(\psi\), the statement having no instance.  Thus \(\bottom\) is a Craig
interpolant.

\emph{The converse.}  Suppose \(\varphi\) is satisfiable and pick
\(v\entails\varphi\).  Then \(v\mathbin{\upharpoonright}X\in U_\varphi\) with
\(v\) as witness, so by Clause~2 of
Lemma~\ref{thm:p02-w02-p02-015-output-content}, \(v\entails\theta\).  Since
no valuation satisfies \(\bottom\), the formulas \(\theta\) and \(\bottom\)
take different values at \(v\) and therefore are different strings.
\end{proof}
\end{annexitem}

\begin{annexitem}{annex:p02-w02-p02-015-valid}{thm:p02-w02-p02-015-valid}{The valid endpoint}
\begin{proof}
Since \(\psi\) is valid, every valuation satisfying \(\varphi\) satisfies
\(\psi\), so \(\varphi\entails\psi\) and the method applies.

\emph{Clause 1.}  This is Theorem~\ref{thm:p02-w02-p02-015-inconsistent}.

\emph{Clause 2.}  The hypotheses of Clause~1 of
Theorem~\ref{thm:p02-w02-p02-015-empty-shared} hold: \(\varphi\entails\psi\),
\(X=\varnothing\), and \(\varphi\) is satisfiable.  That clause gives
\(\theta=\topv\).

\emph{Clause 3.}  Pick \(v\entails\varphi\); then
\(v\mathbin{\upharpoonright}X\in U_\varphi\), so \(m\geq 1\) and \(\theta\)
is the lexicographically ordered disjunction of
Definition~\ref{def:p02-w02-p02-015-output}.  By the proof of Clause~3 of
Lemma~\ref{thm:p02-w02-p02-015-output-content}, each
\(\operatorname{Var}(\kappa_{u_j})=X\) when \(k\geq 1\), and the union over
\(j=1,\ldots,m\) with \(m\geq 1\) is \(X\).

\emph{Clause 4.}  Let \(\varphi\) be valid and \(\varphi\entails\psi\).
Every valuation satisfies \(\varphi\) and hence \(\psi\), so \(\psi\) is valid
and the clause is a special case of the theorem's hypothesis.  Given any
\(u:X\to\{0,1\}\), define \(v_u(p)=u(p)\) for \(p\in X\) and
\(v_u(p)=\mathbf{F}\) otherwise.  Then \(v_u\entails\varphi\) by validity and
\(v_u\mathbin{\upharpoonright}X=u\), so \(u\in U_\varphi\).  Hence
\(U_\varphi\) is the set of all assignments \(X\to\{0,1\}\), which has
\(2^{k}\) members.  For any valuation \(v\),
\(v\mathbin{\upharpoonright}X\in U_\varphi\), so \(v\entails\theta\) by Clause~2
of Lemma~\ref{thm:p02-w02-p02-015-output-content}; thus \(\theta\) is valid.
If \(X=\varnothing\), then \(\varphi\), being valid, is satisfied by the
valuation sending every variable to \(\mathbf{F}\), and Clause~2 gives
\(\theta=\topv\).  If \(X\neq\varnothing\), Clause~3 gives
\(\operatorname{Var}(\theta)=X\neq\varnothing\), while
\(\operatorname{Var}(\topv)=\varnothing\); two strings with different
occurrence sets are different, so \(\theta\neq\topv\).

\emph{Clause 5.}  \(\varphi\entails\theta\): if \(v\entails\varphi\), then
\(v\mathbin{\upharpoonright}X\in U_\varphi\) with \(v\) as witness, and Clause~2
of Lemma~\ref{thm:p02-w02-p02-015-output-content} gives \(v\entails\theta\).
\(\theta\entails\psi\): every valuation satisfies \(\psi\), in particular
every valuation satisfying \(\theta\).  \(\operatorname{Var}(\theta)\subseteq
X\) is Clause~3 of Lemma~\ref{thm:p02-w02-p02-015-output-content}.

For \(\topv\): every valuation satisfies \(\topv\)
(Lemma~\ref{thm:p02-w02-p02-015-variable-free}), so \(\varphi\entails\topv\);
every valuation satisfies \(\psi\), so \(\topv\entails\psi\); and
\(\operatorname{Var}(\topv)=\varnothing\subseteq X\).  Hence \(\topv\) meets
the three conditions of Definition~\ref{def:p02-interpolant}.
\end{proof}
\end{annexitem}
